\documentclass[letterpaper]{article}
\usepackage{iccc}
\usepackage[T1]{fontenc}
\usepackage{graphicx}
\usepackage{booktabs}
\usepackage{url}
\usepackage[hidelinks]{hyperref}
\pdfinfo{
/Title (Small Weather: A Grammar-Based Generator of Eight-Line Weather Poems)
/Author (Ahmad Elhage)
}

\newcommand{\code}[1]{\texttt{\small #1}}

\title{Small Weather: A Grammar-Based Generator\\of Eight-Line Weather Poems}
\author{Ahmad El Hage \textbullet\ Emrecan Birol\\
Course: Computational Creativity\\
Project Assignment 2: Bad Poets Society\\
a.el.hage@umail.leidenuniv.nl \textbullet\ e.birol@umail.leidenuniv.nl}

\begin{document}
\maketitle

\begin{abstract}
\emph{Small Weather} is a small poem generator written in Python. It
writes eight-line free-verse poems about a rainy city, a misty harbour
or a frosted garden. In each poem an ordinary object takes on an
emotional trace and comes back near the end. The system expands a
hand-written grammar, makes 48 seeded candidate poems and picks one
with a simple craft heuristic. Three of its poems were illustrated with
a text-to-image model and typeset as a short anthology. This report
describes the inspiring set, the design of the generator and its
stopping criteria, a mechanical audit of 600 outputs and how the
generative-AI presentation was made. It closes with a reflection on how
the creativity of a system like this could, and could not, be
evaluated. The main finding is that our selection step raises the
heuristic score but makes the output less diverse. This shows how
easily an internal measure of quality can be confused with creativity.
\end{abstract}

\section{Introduction}
Poetry generation is one of the oldest problems in computational
creativity. Approaches range from templates and grammars to evolutionary
search, corpus-based methods and large language models
\citep{lamb2017taxonomy}. For this assignment we wanted a system whose
behaviour we could explain completely, line by line. So we chose a
\emph{finite, parameterised grammar} with a bounded generate-and-select
loop rather than a pretrained model. That choice limits what the system
can do, and much of this report is about those limits.

We worked in five steps. (1)~We gathered a small inspiring set and wrote
down what we wanted to take from each poem. (2)~We designed a fixed
poetic form and turned it into a grammar. (3)~We implemented and tested
the generator, including its stopping and selection rules. (4)~We
generated three poems and presented them with AI illustrations in an
anthology. (5)~We audited the outputs and thought about how
creativity might be evaluated. The project log (\code{LOGBOOK.md})
records each step in more detail.

\section{The Inspiring Set}
\citet{ritchie2007} defines the \emph{inspiring set} as the artefacts
that guide the construction of a creative system, whether or not the
system reads them. Ours is deliberately small and narrow. It has three
short English weather poems from public-domain collections, which we
read on Poets.org and Lit2Go:
\begin{itemize}\itemsep 0pt
\item \textbf{Carl Sandburg, \emph{Fog}} \citep{sandburg1916}. A brief
weather image becomes an animate presence in a concrete city setting.
\emph{Take:} brevity, and weather that \emph{acts}. \emph{Avoid:}
reusing its central animal metaphor.
\item \textbf{Emily Dickinson, \emph{Beclouded}} \citep{dickinson1890}.
The sky and wind behave like people with moods, and a tiny scene
carries a wider emotional implication. \emph{Take:} human verbs for
non-human things. \emph{Avoid:} imitating her rhyme and dashes.
\item \textbf{Sara Teasdale, \emph{There Will Come Soft Rains}}
\citep{teasdale1920}. Gentle sensory details lead to a turn in
perspective about human absence. \emph{Take:} moving from observation
to a turn, and ending on an image that continues beyond the speaker.
\emph{Avoid:} its war setting and rhymed couplets.
\end{itemize}
We picked these poems because they share a scale, which is \emph{small
weather}: one place, one moment and an emotion that is implied rather
than stated. A larger, mixed corpus would have pulled the design in too
many directions. The program does not load these texts, and no source
line appears in its data. The phrase banks are original (see the
\emph{Use of Generative AI} statement below). We also rejected a Markov
chain trained on the set. With only three short poems it would either
produce incoherent text or copy its sources word for word, the
``plagiarism'' end of the novelty scale that \citet{ritchie2007}
warns about.

\section{The Generator}
\subsection{Form}
Our initial idea, recorded before implementation, was a fixed form of
four couplets (eight lines) in free verse. Each line has a set rhetorical
function: (1)~weather acts on a place; (2)~an object holds a memory;
(3)~a sound; (4)~a visual detail; (5)~a first-person \emph{turn};
(6)~the environment answers with a counterpoint; (7)~the object
\emph{returns}; (8)~a closing image that continues beyond the speaker.
Each line maps to something we took from the inspiring set. Lines 1--2
come from Sandburg and Dickinson, and lines 5--8 from Teasdale.

\subsection{Grammar}
The grammar lives in \code{data/grammar.json}, so it can be read and
edited without touching the code. It has eight template banks, one per
line, with one or two surface variants each. For example, line~2 is
either ``\code{The \{adjective\} \{motif\} \{action\} \{memory\}.}'' or
``\code{A \{adjective\} \{motif\} still \{action\} \{memory\}.}''.
The templates draw on two kinds of phrase bank. \emph{Common} banks
(memories, actions, turns, return actions and titles) are shared by all
poems. \emph{Theme} banks hold four places, four weather actions, four
sounds, four details, four counterpoints, four endings and four motifs
for each theme. The themes are rain (city, waiting), sea mist (harbour,
release) and frost (garden, return).

Two design decisions do most of the work for coherence. First, each
motif is a fixed noun--adjective pair, such as \emph{sodden ticket} or
\emph{crooked gate}, so adjectives always fit their nouns. Second, the
motif is chosen \emph{once per poem}. The same object therefore appears
in lines~2 and~7, which gives every poem a small arc of departure and
return. All motif nouns are singular and all opening verbs are
third-person singular, so agreement errors cannot occur. Each theme
allows about $7.9\times10^{8}$ different poem bodies.

\subsection{Generation, Selection and Stopping}
Given a theme and an integer seed, the generator creates a local
\code{random.Random(seed)}. It then expands the grammar into a fixed
budget of candidate poems (48 by default, at most 1,000). A candidate is
\emph{complete} when it has eight lines. This is our first stopping
criterion: a poem has a fixed length, not an open-ended stream. Each
candidate goes through a mechanical validator. It must have no
duplicate lines, 4--14 words per line, no unresolved placeholders, and
the motif in lines~2 and~7. Valid candidates are then scored by a small
craft heuristic:
\[
s = 100 - 1.25\,L - 3\,R + 0.75\min(A,4).
\]
Here $L$ is the total distance from seven words per line, $R$ counts
repeated content words outside the motif, and $A$ counts lines with
adjacent alliteration. The second stopping criterion is the budget.
Search always ends after exactly the requested number of candidates.
The highest-scoring valid candidate wins, and ties go to the earliest
one. If no candidate is valid, the program stops with an error instead
of looping. The output is one poem per requested seed, as text or as
JSON with full metadata (seed, winning candidate index, score and a
SHA-256 hash of the grammar), so every poem can be reproduced exactly.

\subsection{Implementation}
The generator (\code{poet.py}, about 230 lines) uses only the Python
standard library and runs offline. A command-line interface exposes
the theme, seed, poem count, candidate budget, output format and an
alternative grammar file. Six unit tests check determinism, variation
across seeds, bounded failure, argument validation and structural rules
across 150 theme/seed combinations. We also considered Tracery
\citep{compton2015tracery}. A plain Python grammar made it easier to
add a motif that persists across the poem and the validation and
selection steps.

\section{Results}
Figure~\ref{fig:anthology} shows the three anthology poems. They come
from consecutive seeds 42--44 with the default budget, and the winners
were candidates 28, 18 and 23. We did not search for good seeds, and
we did not edit the texts by hand.

\begin{figure*}[t]
\centering
\begin{minipage}[t]{0.32\textwidth}\centering
\includegraphics[width=\linewidth]{01_rain.jpg}\\[3pt]
\begin{minipage}{0.97\linewidth}\scriptsize\raggedright
\textbf{Rain / ticket}\\[2pt]
Under the railway bridge, rain unthreads the evening.\\
The sodden ticket shelters a name nobody calls.\\[2pt]
A loose sign taps against the wall.\\
One lit window floats inside a puddle.\\[2pt]
I thought that leaving would be louder.\\
But the street keeps making room for footsteps.\\[2pt]
The ticket stays when I turn away.\\
A little water travels on without me.
\end{minipage}
\end{minipage}\hfill
\begin{minipage}[t]{0.32\textwidth}\centering
\includegraphics[width=\linewidth]{02_tide.jpg}\\[3pt]
\begin{minipage}{0.97\linewidth}\scriptsize\raggedright
\textbf{A place for sea mist}\\[2pt]
Below the empty pier, sea mist unbuttons the horizon.\\
The paint-worn boat keeps a name nobody calls.\\[2pt]
A chain knocks softly on the stone.\\
All this time, the tide erases one more waterline.\\[2pt]
Perhaps this is how a place forgives.\\
And yet a boat turns slowly towards the open water.\\[2pt]
Even now, the boat has learned to wait.\\
The shore remains, though nothing here stands still.
\end{minipage}
\end{minipage}\hfill
\begin{minipage}[t]{0.32\textwidth}\centering
\includegraphics[width=\linewidth]{03_frost.jpg}\\[3pt]
\begin{minipage}{0.97\linewidth}\scriptsize\raggedright
\textbf{Frost / gate}\\[2pt]
Frost practises its careful handwriting beneath the apple branches.\\
The crooked gate shelters a small unfinished song.\\[2pt]
A branch taps once against the greenhouse.\\
The soil keeps all its letters underground.\\[2pt]
I thought that leaving would be louder.\\
But the light returns without an explanation.\\[2pt]
The gate has learned to wait.\\
Under the path, a root goes on deciding.
\end{minipage}
\end{minipage}
\caption{The \emph{Small Weather} anthology: three generated poems
(seeds 42, 43, 44; unedited) with their AI-generated illustrations.}
\label{fig:anthology}
\end{figure*}

Reading the poems critically, we find that the fixed form gives each
one a clear shape. The object returns, the turn lands in line~5 and the
ending points outward. Some combinations work better than we planned.
In \emph{Frost / gate}, the frost's ``careful handwriting'', the
soil's ``letters underground'' and the ``deciding'' root form a
connected metaphor of writing and growth that no single rule produced.
Other combinations are weaker. In the harbour poem the boat ``turns
slowly towards the open water'' and also ``has learned to wait''. A
reader may find this tension evocative or simply inconsistent, and the
system cannot tell the difference. Repetition also appears
\emph{across} the poems: two of the three use ``I thought that leaving
would be louder'', and two use ``a name nobody calls''. The validator
only prevents duplicate lines within a single poem.

\paragraph{Mechanical audit.}
We generated 100 poems per theme (seeds 0--99) with budgets of 1 and 48
candidates, 600 poems in total (Table~\ref{tab:audit}). All 600 passed
the mechanical checks, and every condition produced 100 distinct poem
bodies. A further sample of 20,000 raw candidates per theme found
\emph{no} validator failures. The grammar meets its constraints by
construction, so in practice the ``generate-and-test'' loop is pure
ranking. The validator is only a safety net for edited grammars.

\begin{table}[h]
\centering\small
\begin{tabular}{lcccc}
\toprule
 & \multicolumn{2}{c}{Mean score} & \multicolumn{2}{c}{Distinct lines /800}\\
Theme & 1 cand. & 48 cand. & 1 cand. & 48 cand.\\
\midrule
Rain  & 87.62 & 94.81 & 180 & 154\\
Tide  & 81.33 & 90.06 & 180 & 167\\
Frost & 80.28 & 89.51 & 180 & 163\\
\bottomrule
\end{tabular}
\caption{Audit of 100 poems per condition.}
\label{tab:audit}
\end{table}

Selection raised the heuristic score in every theme, but it also reduced
the number of distinct lines. The score prefers seven-word lines, so it
systematically favours short phrases. Across the 300 selected poems,
``an unfinished goodbye'' fills the memory slot 118 times, against an
expected 50 under uniform choice, and ``has learned to wait'' is the
most common return action. The selector is not improving the poems in
general. It is narrowing them toward one dimension of what we think a
good line looks like. This trade-off between value and variety is
discussed below.

\section{Presenting the Poems with Generative AI}
We presented the poems as a small illustrated anthology
(\code{Small\_Weather\_Anthology.pdf}), with one page per poem. This
format suits short, quiet poems and keeps the exact text easy to read.
For each poem we gave the built-in image generation tool of our AI
assistant a structured prompt. The prompt contained the \emph{complete
poem} as a semantic reference, a scene description centred on the
poem's object and setting (a ticket under a railway bridge, a boat
below a pier, a gate by a greenhouse) and a shared style: gouache and
fine ink on ivory paper, in a muted slate, teal and ochre palette. It
also listed constraints: no lettering, no people and no objects turned
into cartoon characters. The first image from each prompt was kept
without editing. A ReportLab script then set each illustration above
its poem, so the text is real, selectable type. We did not ask the
image model to render the words, so it could not misspell or rewrite
the poem. All prompts, file hashes and tool details are stored in
\code{presentation/provenance.json}. The provider did not expose a
model identifier or image seed, so the images cannot be regenerated
exactly.

\paragraph{Reflection.}
The image model was most useful for unifying the collection. A shared
style made three separately generated poems look like a deliberate
series, and that probably increases how much readers value them
\citep{colton2008}. The images also \emph{interpret}. The rain image
makes the ticket and the reflected window concrete and gives the
departure a mood. The frost image shows exposed roots and turns the
poem's ambiguous ``deciding'' root into something literal. The harbour
image captures stillness and tethering but cannot show
``forgiveness''. An illustration can therefore narrow a reader's
imagination as easily as it widens it. We also noticed that the
scene-and-style text we wrote shaped the images more than the poem
did. That is our creative contribution, not the generator's. It also
raises the question of where the credit lies when a poem's perceived
value rises because of a picture the poet did not make.

\section{Evaluating the Creativity of the System}
We did not run a formal creativity study, so this section is a
reflection. It draws on the audit and on established frameworks.

\paragraph{What the numbers do not show.}
The audit establishes reproducibility, bounded termination and
structural correctness. The fact that 100 out of 100 poem bodies are
distinct says nothing about creativity, because a grammar with
$10^{8}$ combinations will almost never repeat itself in a small
sample. The heuristic score is even less informative. It measures
exactly what the selector optimises, so it is not independent
evidence. In fact the score went up while diversity went down.

\paragraph{Ritchie's criteria.}
\citet{ritchie2007} proposes judging a system by the \emph{typicality}
and \emph{quality} of its outputs and by their novelty relative to the
inspiring set. Our system is probably strong on typicality, since every
output is recognisably the intended form. It may be moderate on
quality, which only human ratings could estimate. It is weak on novelty
\emph{within} its own space. Every line comes from a fixed bank, and
the system can never produce a phrase we did not write. Novelty
relative to the inspiring set is high only because we made it so by
design.

\paragraph{Boden and the creative tripod.}
In \citeauthor{boden2004}'s terms \citeyearpar{boden2004}, the system
is \emph{combinational} and at best weakly \emph{exploratory}. It
searches a conceptual space we defined and cannot change that space.
Following \citet{wiggins2006}, it can traverse its rules but has no
mechanism to transform them. Against \citeauthor{colton2008}'s
creative tripod \citeyearpar{colton2008}, it shows some \emph{skill} in
producing well-formed, coherent poems and very limited
\emph{appreciation} in its craft score. It shows almost no
\emph{imagination}, since it never departs from its templates.
It also does not frame its work: it cannot explain why a poem matters
\citep{colton2012fullface}.

\paragraph{How we would evaluate it.}
A meaningful evaluation would ask human readers. We would run a blind
study with a pre-registered seed schedule. Participants would rate
poems from the 1-candidate and 48-candidate conditions for coherence,
imagery, surprise and emotional effect. A second, counterbalanced
comparison of plain-text and illustrated versions would test whether
the presentation changes perceived value. Following
\citet{jordanous2012}, we would first state which components of
creativity matter in this domain, such as surprise, value and
intentionality, and evaluate against those rather than against a
single score. Finally, a ``mere generation'' baseline (one candidate,
no selection) is essential. Our audit suggests that a critic can make
poems \emph{more polished and more alike} at the same time, and only
readers can say whether that is a gain.

\section{Limitations and Future Work}
The structure, themes and almost all of the phrasing were fixed by
humans with AI assistance, so most of the creative decisions sit
outside the running program. The most valuable next step would be to
give the system ways to change itself. For example, it could penalise
phrases reused across a collection, learn new motifs from a larger
inspiring set or vary the form. A second step would be to have it
write a short framing statement for each poem.

\section{Use of Generative AI}
Generative AI was used at three distinct stages. (1)~\emph{Development:}
an AI coding assistant (OpenAI Codex) helped choose the concept, draft
the code and the original phrase banks, check outputs and write
documentation. The author reviewed and directed this work, and an AI
assistant (Claude) helped draft this report. (2)~\emph{Runtime:} the
poem generator itself uses no model or network. It only recombines the
saved grammar. (3)~\emph{Presentation:} the three illustrations were
produced by a text-to-image tool from the prompts described above. The
poem texts were never altered by any AI tool after generation.


\begin{thebibliography}{99}

\bibitem[Boden(2004)]{boden2004}
Boden, M.~A. 2004.
\newblock \emph{The Creative Mind: Myths and Mechanisms}.
\newblock London: Routledge, 2nd edition.

\bibitem[Colton, Goodwin, and Veale(2012)]{colton2012fullface}
Colton, S.; Goodwin, J.; and Veale, T. 2012.
\newblock Full-{FACE} poetry generation.
\newblock In \emph{Proceedings of the Third International Conference on
Computational Creativity}, 95--102.

\bibitem[Colton(2008)]{colton2008}
Colton, S. 2008.
\newblock Creativity versus the perception of creativity in
computational systems.
\newblock In \emph{AAAI Spring Symposium: Creative Intelligent Systems},
14--20.

\bibitem[Compton, Kybartas, and Mateas(2015)]{compton2015tracery}
Compton, K.; Kybartas, B.; and Mateas, M. 2015.
\newblock Tracery: An author-focused generative text tool.
\newblock In \emph{Interactive Storytelling (ICIDS 2015)}, 154--161.
Springer.

\bibitem[Dickinson(1890)]{dickinson1890}
Dickinson, E. 1890.
\newblock Beclouded.
\newblock In \emph{Poems by Emily Dickinson, Series One}. Boston: Roberts
Brothers. Read via Lit2Go, University of South Florida.
\newblock \url{https://etc.usf.edu/lit2go/114/the-poems-of-emily-dickinson-series-one/2632/nature-poem-29-beclouded/}.

\bibitem[Jordanous(2012)]{jordanous2012}
Jordanous, A. 2012.
\newblock A standardised procedure for evaluating creative systems:
Computational creativity evaluation based on what it is to be creative.
\newblock \emph{Cognitive Computation} 4(3):246--279.

\bibitem[Lamb, Brown, and Clarke(2017)]{lamb2017taxonomy}
Lamb, C.; Brown, D.~G.; and Clarke, C.~L.~A. 2017.
\newblock A taxonomy of generative poetry techniques.
\newblock \emph{Journal of Mathematics and the Arts} 11(3):159--179.

\bibitem[Ritchie(2007)]{ritchie2007}
Ritchie, G. 2007.
\newblock Some empirical criteria for attributing creativity to a
computer program.
\newblock \emph{Minds and Machines} 17(1):67--99.

\bibitem[Sandburg(1916)]{sandburg1916}
Sandburg, C. 1916.
\newblock Fog.
\newblock In \emph{Chicago Poems}. New York: Henry Holt. Read via
Poets.org, \url{https://poets.org/poem/fog}.

\bibitem[Teasdale(1920)]{teasdale1920}
Teasdale, S. 1920.
\newblock There will come soft rains.
\newblock In \emph{Flame and Shadow}. New York: Macmillan. Read via
Poets.org, \url{https://poets.org/poem/there-will-come-soft-rains}.

\bibitem[Wiggins(2006)]{wiggins2006}
Wiggins, G.~A. 2006.
\newblock A preliminary framework for description, analysis and
comparison of creative systems.
\newblock \emph{Knowledge-Based Systems} 19(7):449--458.

\end{thebibliography}
\end{document}